\documentclass[a4paper]{article}

\usepackage[utf8]{inputenc}
\usepackage[T1]{fontenc}
\usepackage[ngerman]{babel}
\usepackage{xcolor}

\begin{document}
    
    \section{Der Infinitivsatz}
        
        Ein Infinitivsatz ist ein Nebensatz mit \textbf{zu + Infinitiv}.
            
            {\small\textcolor{gray}{An infinitive clause is a subordinate clause with \textbf{zu + infinitive}.}}
        
        Beispiel:
        
        \textbf{Ich versuche, Deutsch zu lernen.}
        
        {\small\textcolor{gray}{I try to learn German.}}
        
        Hier ist \textbf{zu lernen} der Infinitivsatz.
            
            {\small\textcolor{gray}{Here, \textbf{zu lernen} is the infinitive clause.}}
        
        Der Infinitivsatz steht meistens nach einem Komma.
            
            {\small\textcolor{gray}{The infinitive clause usually comes after a comma.}}
        
        Beispiel:
        
        \textbf{Ich habe keine Zeit, heute zu lernen.}
        
        {\small\textcolor{gray}{I have no time to study today.}}
        
        Bei trennbaren Verben kommt \textbf{zu} in die Mitte des Verbs.
            
            {\small\textcolor{gray}{With separable verbs, \textbf{zu} goes into the middle of the verb.}}
        
        Beispiele:
        
        \textbf{anfangen $\rightarrow$ anzufangen}
        
        {\small\textcolor{gray}{to begin $\rightarrow$ to begin}}
        
        \textbf{aufstehen $\rightarrow$ aufzustehen}
        
        {\small\textcolor{gray}{to get up $\rightarrow$ to get up}}
        
        Beispiel:
        
        \textbf{Ich habe vor, morgen früh aufzustehen.}
        
        {\small\textcolor{gray}{I plan to get up early tomorrow.}}

\end{document}